\par\medskip
\noindent\textbf{Observación 4.16.}
La teoría global del complejo cotangente relativo---reducida aquí al
complejo truncado de longitud uno, que sin embargo basta para muchos
fines de la teoría de deformaciones---generaliza y simplifica
resultados como 4.10. Sea $S$ un esquema, y sean
$S_0\subset S'\subset S$ subesquemas definidos por ideales
$J\supset I$ con $JI=0$. Sea $X'$ plano sobre $S'$, y escribamos
\[
L_\bullet=L^{X_0/S_0}
\]
para el complejo cotangente relativo de $X_0$ sobre $S_0$,
considerado como un objeto de $D(X_0,\mathcal O_{X_0})$. Sea
$\underline F$ la categoría de los $S$-esquemas planos $X$ que
extienden $X'/S'$; así, un objeto está provisto de un
$S'$-isomorfismo
\[
X\times_SS'\xrightarrow{\ \sim\ }X'.
\]
Entonces:

\smallskip
\noindent\textbf{(a)} Para un objeto $X$ de $\underline F$, su grupo
de automorfismos que inducen la identidad sobre $X'$ es canónicamente
isomorfo a
\[
\operatorname{Ext}^0(X_0;L_\bullet,I_{X_0}),
\]
donde $I_{X_0}$ es la imagen inversa sobre $X_0$ del
$\mathcal O_{S_0}$-módulo $I$.

%%%%%%%%%%  scan idx 123  |  folio 112  %%%%%%%%%%

\smallskip
\noindent\textbf{(b)} El conjunto de clases de isomorfía de objetos
de $\underline F$ es vacío o bien un torsor (un espacio homogéneo
principal) bajo $\operatorname{Ext}^1(X_0;L_\bullet,I_{X_0})$.

\smallskip
\noindent\textbf{(c)} Existe un elemento canónico
\[
c\in\operatorname{Ext}^2(X_0;L_\bullet,I_{X_0})
\]
cuya anulación es necesaria y suficiente para que $\underline F$ no
sea vacía.

Para estudiar estos grupos, a menudo conviene escribir
\[
\operatorname{Ext}^i(X_0;L_\bullet,I_{X_0})
=H^i\bigl(X_0,R\underline{\operatorname{Hom}}
  (L_\bullet,I_{X_0})\bigr).
\]
Los haces de cohomología de
$R\underline{\operatorname{Hom}}(L_\bullet,I_{X_0})$ son
precisamente los haces $\underline D^i_{X_0}(I_{X_0})$ introducidos
en la sección 3. Para $i=0,1,2$, pueden considerarse,
respectivamente, como el haz de automorfismos de un objeto $X$ de
$\underline F$ (automorfismos que inducen la identidad sobre $X'$),
el haz que mide la indeterminación de las clases de isomorfía de las
soluciones locales del problema de deformación plana y el haz de
obstrucciones a la existencia de soluciones locales.

La sucesión espectral estándar es
\[
E_2^{p,q}
=H^p\bigl(X_0,\underline D^q_{X_0}(I_{X_0})\bigr)
\Longrightarrow
\operatorname{Ext}^{p+q}(X_0;L_\bullet,I_{X_0}).
\]
Proporciona los isomorfismos conocidos
\begin{align*}
\operatorname{Aut}(X/X')
&\simeq H^0\bigl(X_0,\underline D^0_{X_0}(I_{X_0})\bigr)\\
&\simeq H^0\bigl(X_0,
  \underline{\operatorname{Hom}}
  (\Omega^1_{X_0/S_0},I_{X_0})\bigr)\\
&=\operatorname{Hom}(\Omega^1_{X_0/S_0},I_{X_0}),
\end{align*}
y la sucesión exacta de grados bajos
\[
0\longrightarrow H^1(X_0,\underline D^0_{X_0}(I_{X_0}))
\longrightarrow\operatorname{Ext}^1
\longrightarrow H^0(X_0,\underline D^1_{X_0}(I_{X_0}))
\longrightarrow H^2(X_0,\underline D^0_{X_0}(I_{X_0}))
\longrightarrow\operatorname{Ext}^2.
\]
En particular, si
$H^2(X_0,\underline D^0_{X_0}(I_{X_0}))=0$ y existen deformaciones
planas locales de $X'/S'$, entonces las deformaciones locales
compatibles se pegan en una deformación plana global. Más
generalmente, para una sección prescrita del haz de clases de
isomorfía de deformaciones planas locales, el argumento de 4.10 sitúa
la obstrucción a una deformación global compatible en
$H^2(X_0,\underline D^0_{X_0}(I_{X_0}))$.

%%%%%%%%%%  scan idx 124  |  folio 113  %%%%%%%%%%

La misma sucesión espectral muestra que $\operatorname{Ext}^2$ se
anula siempre que
\begin{equation}\tag{*}
H^0(X_0,\underline D^2_{X_0})=0,
\qquad
H^1(X_0,\underline D^1_{X_0})=0,
\qquad
H^2(X_0,\underline D^0_{X_0})=0.
\end{equation}
La primera igualdad se cumple siempre cuando $X_0$ es una
intersección completa relativa sobre $S_0$, pues entonces
$\underline D^2_{X_0}(M)=0$ para todo $\mathcal O_{X_0}$-módulo $M$.
La segunda se cumple cuando el soporte de $\underline D^1_{X_0}$ es
discreto, como ocurre cuando $X_0$ es liso sobre $S_0$ fuera de un
subconjunto discreto. Si se cumplen las tres igualdades de $(*)$,
entonces $\underline F$ no es vacía por (c), lo que generaliza
4.13(b).

Existe también una interpretación geométrica que solo requiere los
funtores globalizados $\underline D^i_{X_0}$, en lugar de la teoría
global del complejo cotangente relativo. Las tres obstrucciones
sucesivas a extender planamente $X'/S'$ sobre $S$ pertenecen a los
tres grupos mostrados en $(*)$. La teoría local define la primera
obstrucción
\[
c_1\in H^0(X_0,\underline D^2_{X_0})
\]
como la obstrucción a la existencia de soluciones locales. Si
$c_1=0$, el haz $\underline F$ de clases de isomorfía de soluciones
locales es canónicamente un torsor bajo $\underline D^1_{X_0}$. Su
obstrucción a tener una sección es
\[
c_2\in H^1(X_0,\underline D^1_{X_0});
\]
tal sección es un sistema coherente de clases de isomorfía de
soluciones locales. Si $c_2=0$, elijamos una sección de
$\underline F$. La obstrucción a una solución global que realice las
clases locales elegidas es entonces
\[
c_3\in H^2(X_0,\underline D^0_{X_0}).
\]
Cambiar la sección de $\underline F$ por una sección de
$\underline D^1_{X_0}$ da una aplicación
\[
d:H^0(X_0,\underline D^1_{X_0})
\longrightarrow H^2(X_0,\underline D^0_{X_0}).
\]
Existe una solución si y solo si $c_1=0$, $c_2=0$ y la clase de
$c_3$ en $\operatorname{Coker}d$ es cero (estando definido cada
$c_i$ únicamente después de que se anulen las obstrucciones
precedentes).

%%%%%%%%%%  scan idx 125  |  folio 114  %%%%%%%%%%

Los tres grupos de obstrucciones
\[
H^0(X_0,\underline D^2_{X_0}),
\qquad
H^1(X_0,\underline D^1_{X_0}),
\qquad
H^2(X_0,\underline D^0_{X_0})/\operatorname{Im}d
\]
son los términos iniciales $E_2^{0,2}$, $E_2^{1,1}$ y $E_3^{2,0}$
de la sucesión espectral precedente; salvo el signo, $d$ es su
diferencial
\[
d_2:E_2^{0,1}\longrightarrow E_2^{2,0}.
\]
